% Candidate A.54. Unregistered pending independent parent proof/source audit.
\subsection{The degree-one-hundred-twenty-five Heisenberg cover}

\begin{statement}{A.54}{Proved}{An actual non-CM factor and the nonsplit-prime normalizers}
Let \(C\) be the smooth proper curve
\[
 \begin{aligned}
 y^2&=x^3-x,\qquad z^5=\psi _5(x),\\
 \psi _5(x)&=5x^{12}-62x^{10}-105x^8+300x^6-125x^4+50x^2+1,
 \end{aligned}
                                                               \tag{A.54.1}
\]
and let \(q:C\to E\), \(S=q^{-1}(E[5])\), and \(X=C\setminus S\).
The map \([5]q:X\to U\) is a geometric Galois cover
with Heisenberg group of order \(125\), genus \(51\), and \(25\) cusps.
After a finite number-field extension,
\[
 J_C\sim E\times A_2^5\times B_8^5,\qquad
 A_2\sim J(v^2=w^5-1),\quad \dim B_8=8.             \tag{A.54.2}
\]
Here \(A_2\) has primitive quartic CM by \(\mathbf Q(\zeta _5)\),
and \(B_8\) has an action by \(K=\mathbf Q(\zeta _{20})\)
but has no geometric CM abelian subvariety.
The last assertion follows from an actual genus-eleven quotient of \(C\)
and two persistent, unequal quadratic Frobenius extensions of \(K\).

For every coefficient prime \(\ell\not\equiv1\pmod5\),
the nontrivial procyclic subgroups of the geometric pro-\(\ell\) group
with open arithmetic normalizer are exactly the nontrivial subgroups
of cusp inertia conjugates. The full pro-five fixed subgroup for every
packet and open derived subgroup in Part A is the selected cusp inertia.
This statement leaves the normalizer question for \(\ell\equiv1\pmod5\),
and the other full fixed-subgroup questions, open.
\end{statement}

\begin{proof}
\paragraph{The actual cover.}
The division-polynomial divisor is
\[
 \operatorname{div}(\psi _5)=D_{E[5]}-25(O),
 \qquad D_T:=\sum_{P\in T}(P).                     \tag{A.54.3}
\]
Its orders are \(1\) at the twenty-four nonzero points and \(-24\)
at \(O\). The Kummer cover is connected, of degree five and totally
ramified at exactly these twenty-five points. Consequently
\(2g(C)-2=25(5-1)=100\); removing their unique preimages gives
the asserted degree and number of cusps.

Every translation by \(T\in E[5]\) lifts after a finite constant
extension. Indeed, the divisor of \(\psi _5(P+T)/\psi _5(P)\) is
\(25((O)-(-T))\). Since \(5((O)-(-T))\) is principal, this
ratio is a constant times a fifth power of a rational function.
Taking its fifth root gives the lift. The lifts and
\(c:z\mapsto\zeta _5z\) form a central extension of \(E[5]\)
by \(\mu _5\), of order \(125\). It is nonabelian: in an abelian
cover of a once-punctured elliptic curve, the sole peripheral
commutator has trivial image, whereas here its inertia has order five.

This group has exponent five. In a class-two group with central
kernel of order five, the fifth-power map on lifts defines an additive
map \(E[5]\to\mu _5\). The actual involution
\(\sigma(x,y,z)=(x,-y,z)\) acts by \(-1\) on \(E[5]\) and
trivially on the central kernel. Equivariance therefore makes that
additive map equal to its negative, hence zero. The nonabelian
exponent-five central extension is the Heisenberg group \(H_5\).
All lifts, constants, and later projectors are defined over a finite
number field. A.22 applies to this actual five-group cover, giving
the asserted full pro-five fixed subgroup.

\paragraph{The multiplicity factors are algebraic factors of this curve.}
Put \(F=\mathbf Q(\zeta _5)\). The nontrivial-central component
of \(\mathbf Q[H_5]\) is \(M_5(F)\): its standard representation
on \(F^5\) is generated by
\[
 Ue_j=e_{j+1},\qquad Ve_j=\zeta _5^j e_j
                         \quad(j\bmod5).
\]
The central-invariant homology is \(H_1(E,\mathbf Q)\); all
translations act trivially there. The remaining rational dimension
is \(102-2=100\), or five copies of the simple rational module
of dimension twenty. Matrix idempotents, cleared of denominators,
thus give actual isogeny factors \(J_C\sim E\times B_{10}^5\),
where \(\dim B_{10}=10\) and \(\dim_F H_1(B_{10})=5\).

The actual order-four automorphism is
\(\tau(x,y,z)=(-x,iy,z)\), with \(\tau^2=\sigma\).
Choose a basis \(T,iT\) of \(E[5]\), avoiding its two
\(i\)-eigenlines. A central adjustment of a lift of \(T\)
ensures that \(\sigma\) takes that lift to its inverse:
the adjustment solves one linear equation with coefficient two
in \(\mathbf F_5\). Taking its \(\tau\)-conjugate for the
second lift gives \(U\mapsto V,V\mapsto U^{-1}\), with a
consistent choice of central generator. The Fourier intertwiner
is
\[
 A=\frac1{\sqrt5}(\zeta _5^{jk})_{0\le j,k\le4},
 \qquad \sqrt5=1+2\zeta _5+2\zeta _5^4.
                                                               \tag{A.54.4}
\]
Direct multiplication gives \(AUA^{-1}=V\),
\(AVA^{-1}=U^{-1}\), and \(A^2e_j=e_{-j}\).
Thus \(A^4=1\), and its multiplicities at \(1,-1,i,-i\)
are \((2,1,1,1)\). Write \(\tau=A\otimes T\) on the
noncentral homology. Equivalently, \(A^{-1}\tau\) is an actual
rational endomorphism commuting with \(M_5(F)\) and acting
as \(T\) on \(B_{10}\).

The quotients by \(\tau\) and \(\sigma\) have equations
\(z^5=P_6(s)\), \(s=x^2\), and \(z^5=\psi _5(x)\),
respectively. Their genera are ten and twenty-two: the polynomials
have six and twelve simple roots, with degrees coprime to five.
The original elliptic summand contributes no invariants under either
automorphism. If the multiplicities of \(T\) at \(1,-1,i,-i\)
are \(a,b,d,d\), equality of the last two follows from \(i\notin F\),
and the actual quotient genera give
\[
 a+b+2d=5,\qquad 2a+b+2d=5,
 \qquad3(a+b)+4d=11.
\]
Hence \((a,b,d,d)=(0,1,2,2)\). The \(T=-1\) factor has
dimension two and rank one over \(F\), so is CM; the \(T^2=-1\)
factor has dimension eight, with \(K=F(i)\) acting at rank two.
These are the factors \(A_2,B_8\) in (A.54.2).

\paragraph{An explicit quotient identifies the same factors over a number field.}
Take
\[
 \alpha=2-i,\qquad \lambda_H=\frac{1-2i}{5},\qquad
 \lambda=\frac{1+2i}{5},\qquad h_H=x^2-\lambda_H.
\]
The nonzero points of \(H=\ker\alpha\) have \(x^2=\lambda_H\).
This follows either from the duplication law, with
\([2]P=[i]P\), or from the following exact coordinate identity:
\[
 \begin{aligned}
 x(\alpha P)&=
 \frac{x(3+4i)(x^2-1-2i)^2}{25(x^2-\lambda_H)^2},
 \\
 \psi _5&=(-35-120i)\bigl(x(\alpha P)^2-\lambda\bigr)h_H^5.
 \end{aligned}
                                                               \tag{A.54.5}
\]
Its divisor explanation is equally exact:
\(\operatorname{div}(h_H)=D_H-5(O)\), while
\(x^2-\lambda\) has divisor \(D_{\ker\bar\alpha}-5(O)\).
Pulling the latter back by \(\alpha\) gives
\(D_{E[5]}-5D_H\). The remaining ratio is constant, and
the expansion at \(O\) identifies it as \(5\alpha^4=-35-120i\).

After adjoining \(\gamma\) with \(\gamma^5=-35-120i\),
the actual map
\[
 \begin{gathered}
 C\longrightarrow D_\lambda:\qquad
 (P,z)\longmapsto\left(\alpha P,\frac{z}{\gamma h_H(P)}\right),
 \\
 D_\lambda:\quad y^2=x^3-x,\quad z^5=x^2-\lambda.
 \end{gathered}
                                                               \tag{A.54.6}
\]
has degree five. The identity makes its function-field extension
the pullback along \(\alpha\), so it is the quotient by a
noncentral order-five subgroup, denoted again by \(H\).
That subgroup acts freely since all inertia of the full cover is
central. It is \(\tau\)-stable, and (A.54.6) is equivariant
for \(\tau\) and \(c\). In particular \(D_\lambda\) has
genus eleven. This explicit identity, rather than branch data alone,
identifies the quotient and its constant twist.

In the standard five-dimensional module, the invariant line of
every \(\sigma\)-stable noncentral subgroup has \(A^2=+1\).
For completeness, such a subgroup is generated by
\(\zeta _5^{ab/2}U^aV^b\); if \(a\ne0\), its invariant
vector has coordinates \(\zeta _5^{bj^2/(2a)}\), which are
even in \(j\), and for \(a=0\) it is \(e_0\).
Thus on the noncentral homology of this actual quotient,
\(\sigma=T^2\). It follows that
\[
 J_{D_\lambda}\sim E\times A_2\times B_8,\qquad
 D_\lambda/\sigma:\quad x^2=z^5+\lambda.             \tag{A.54.7}
\]
The last curve becomes \(v^2=w^5-1\) after scaling over
a finite extension, proving the asserted identification of \(A_2\).
Its holomorphic differentials \(dz/x,z\,dz/x\) give the
quartic CM type \(\{1,2\}\subset(\mathbf Z/5)^\times\).

\paragraph{The actual primitive-twenty factor and its Hodge signature.}
On \(D_\lambda\), put \(s=x^2\) and \(W=yz\). Then
\[
 W^{20}=s^5(s-1)^{10}(s-\lambda)^4.                 \tag{A.54.8}
\]
This is a birational cyclic-twenty model. Indeed
\(y^4=s(s-1)^2\), \(z^5=s-\lambda\), and
\(y=W^5/(s(s-1)^2(s-\lambda))\), recovering the two
coprime subextensions. Also \(x=y^2/(s-1)\).
Choose once and for all
\[
 K=\mathbf Q(\zeta),\quad \zeta=\zeta_{20},
 \quad i=\zeta^5,\quad \zeta_5=\zeta^4.
\]
The automorphism \(\delta=\tau c^4\) sends \(W\) to
\(\zeta W\); it satisfies \(\delta^4=c\),
\(\delta^5=\tau\), and \(\delta^{10}=\sigma\).
Consequently its primitive-twenty rational projector selects
exactly the same \(B_8\) as (A.54.7). It defines that factor
over \(K\) itself, independently of the finite extension used
to compare it with \(C\).

The four local exponents in (A.54.8) are \((5,10,4,1)\)
modulo twenty, the last at infinity. The space of holomorphic
differentials with \(\delta\)-eigenvalue \(\zeta^{-k}\)
has dimension
\[
 h_k=-1+\sum_{m\in\{5,10,4,1\}}\left\{\frac{km}{20}\right\}.
                                                               \tag{A.54.9}
\]
One obtains this directly by writing a differential as
\(r(s)\,ds/W^k\) and imposing the four valuation inequalities;
the integral parts give a line bundle on \(\mathbf P^1\)
of degree \(-2+\sum\{km/20\}\), whose space of sections
has the stated dimension. For primitive \(k\), the values are
\[
\begin{array}{c|rrrrrrrr}
k&1&3&7&9&11&13&17&19\\\hline
h_k&0&1&1&1&1&1&1&2.
\end{array}                                                     \tag{A.54.10}
\]
In particular there are balanced embeddings with dimensions
\((1,1)\), and one unbalanced conjugate pair \((0,2)\).
No balanced-determinant or quaternion argument is being substituted
for the following arithmetic calculation.

\paragraph{Exact Frobenius fields at two places of the same field \(K\).}
Both \(41\) and \(61\) split completely in \(K\).
For \(q\) equal to either, choose a primitive root \(g\) modulo
\(q\), and the place with
\(\zeta\mapsto g^{(q-1)/20}\). This fixes
\(i\mapsto g^{(q-1)/4}\) and \(\lambda\mapsto(1+2i)/5\)
in the reduction. The four branch sections in (A.54.8) remain
distinct and all ramification is tame, so the smooth normalization
and its Jacobian have good reduction. Define
\(\chi(g)=\zeta\) and extend \(\chi(0)=0\). On the
primitive eigenspace with \(\delta^*\)-eigenvalue \(\zeta\),
the trace of Frobenius is
\[
 a_q=-\sum_{u\in\mathbf F_q}
       \chi\bigl(u^5(u-1)^{10}(u-\lambda)^4\bigr).  \tag{A.54.11}
\]
This follows from the curve trace formula applied to the
Kummer eigensheaf; every local inertia character is nontrivial,
so compactly supported and proper primitive cohomology agree.
Equivalently, average the ordinary fixed-point formula over
the twenty deck twists. The convention \(\chi(g)=\zeta\)
agrees with the displayed \(\delta^*\) eigenvalue, since
the Frobenius permutation of a fibre is multiplication by
\(f(u)^{(q-1)/20}\).
These are direct applications of
\cite[Theorem 25.1, p.147, and Theorem 29.4, p.165]{PBC_MilneLEC}.

For the square Frobenius use \(\chi\circ N_{\mathbf F_{q^2}/\mathbf F_q}\)
in the same sum. Writing \(u=a+b\sqrt\nu\), each norm is
an elementary quadratic expression, so the retained certificate
uses only integer arithmetic and a specified nonsquare \(\nu\).
Set \(d_q=(a_q^2-a_{q^2})/2\),
\(\Delta_q=a_q^2-4d_q\); the relative Frobenius polynomial is
\(t^2-a_qt+d_q\). The exact sums give
\[
\begin{aligned}
a_{41}&=\zeta^5-\zeta^4+\zeta^3-2\zeta^2+2\zeta-1,\\
a_{61}&=\zeta^7+\zeta^6-2\zeta^5-4\zeta^4
                  +\zeta^3+4\zeta^2-4\zeta-6,\\
\Delta_{41}&=80\zeta^7+60\zeta^6-40\zeta^5-85\zeta^4\\
           &\hspace{1em}+160\zeta^3+75\zeta^2-70\zeta-55,\\
\Delta_{61}&=50\zeta^7+85\zeta^6-170\zeta^5-105\zeta^4\\
           &\hspace{1em}+60\zeta^3+115\zeta^2-170\zeta-30.
\end{aligned}                                                  \tag{A.54.12}
\]
All expressions are reduced modulo
\(\Phi_{20}(\zeta)=\zeta^8-\zeta^6+\zeta^4-\zeta^2+1\).
The complete twenty-bin counts, determinants, and extension-field
parameters are retained with the executed certificate.

The following tiny residue certificates show that the two
quadratic extensions are fields and are unequal as extensions
of the fixed field \(K\):
\[
\begin{array}{c|r|r|r|r}
\text{element}&r&\zeta\bmod r&\text{residue}&
                 \text{quadratic character}\\\hline
\Delta_{41}&61&33&17&-1\\
\Delta_{61}&41&39&15&-1\\
\Delta_{41}&41&39&5&1\\
\Delta_{41}\Delta_{61}&41&39&34&-1
\end{array}                                                     \tag{A.54.13}
\]
Every indicated residue is a unit and every indicated root
satisfies \(\Phi_{20}\). A square in \(K\) which is a unit
at such a place has square residue. Therefore
\[
 L_{41}=K(\sqrt{\Delta_{41}}),\qquad
 L_{61}=K(\sqrt{\Delta_{61}}),\qquad L_{41}\not\simeq_K L_{61}.
                                                               \tag{A.54.14}
\]

These fields survive every positive Frobenius power. For
\(q=41\), at the \(K\)-place \(\zeta\mapsto39\), the
trace is \(13\) modulo \(41\) and the determinant is zero.
For \(q=61\), the corresponding witness is
\(\zeta\mapsto24\), trace \(5\), determinant zero.
The determinant is a nonzero algebraic integer. The Newton
polygon of \(t^2-a_qt+d_q\) at this witness therefore has
one root of valuation zero and one of positive valuation.
Their quotient cannot be a root of unity. Thus, for every
\(n\ge1\), \(K[\pi_q^n]=K[\pi_q]=L_q\).
This is precisely what is needed when the endomorphisms of
\(B_8\) are defined only after finite extension.

\paragraph{Why the field certificates exclude every CM factor.}
First, if \(B_8\) has a nonzero CM isogeny factor, its maximal
CM sub-Hodge structure is preserved by every endomorphism,
hence by \(K\). Its rational dimension is therefore a positive
multiple of eight. A proper such part has abelian dimension four;
its quotient also has dimension four and inherits a faithful
unital \(K\)-action. A rank-one Hodge structure over the CM
field \(K\) is of CM type: its Mumford--Tate group is contained
in \(\operatorname{Res}_{K/\mathbf Q}\mathbf G_m\).
Semisimplicity of polarizable weight-one Hodge structures then
makes the whole \(B_8\) CM. This argument uses the quotient,
and does not assume the existence of a chosen \(K\)-stable
polarization complement.

Suppose, therefore, that \(B_8\) is CM. Put
\[
 \mathcal D=\operatorname{End}_{\mathrm{Hdg},K}
                  (H_1(B_8,\mathbf Q)).
\]
At a balanced embedding in (A.54.10), its complexification
preserves two distinct one-dimensional Hodge summands; hence
\(\dim_K\mathcal D\le2\). The Mumford--Tate group of a CM
Hodge structure is a torus, whose commutant on a two-dimensional
space has dimension at least two. Thus \(\mathcal D\) is a
quadratic etale \(K\)-algebra. These are actual rational
endomorphisms of the abelian variety, by the weight-one full
faithfulness in \cite[Theorems 4.4 and 7.12, pp.489 and 508;
CM discussion on pp.509 and 519]{fghMilneSVModuli}.
Choose a finite extension over which a basis of \(\mathcal D\)
is defined.

At either chosen good place, these endomorphisms specialize
injectively and commute with some positive power \(\pi_q^n\).
Indeed, they extend to the abelian scheme, and prime-to-\(q\)
Tate modules identify generic and special homomorphisms
faithfully. The required integral faithfulness is exactly
\cite[Theorem 10.15, pp.49--50]{fghMilneAV2008}; no surjectivity
part of a Tate isogeny theorem is used here.
In the rational endomorphism algebra of the special fibre consider
\[
 \mathcal C_q=\operatorname{Cent}_{\operatorname{End}^0(\overline B_{8,q})}
                          (K,\pi_q^n).
\]
For an auxiliary prime different from \(q\), its faithful
Tate realization lies in the centralizer of a nonscalar
separable quadratic operator on a rank-two
\(K\otimes\mathbf Q_\ell\)-module. That centralizer has
\(K\)-dimension two. Thus \(\dim_K\mathcal C_q\le2\);
it already contains \(K[\pi_q^n]=L_q\), so
\(\mathcal C_q=L_q\) as a global \(K\)-algebra.
The specialized \(\mathcal D\) has the same dimension and
therefore gives \(\mathcal D\simeq_K L_q\) at both places.
This contradicts (A.54.14).
The residue test at \(r=41\) in (A.54.13) is a number-field
square-class test, not an assertion using a prime-to-characteristic
Tate module at that same prime. This proves that \(B_8\)
has no geometric CM factor.

\paragraph{The normalizers away from split quartic primes.}
After any finite extension defining the factors, the Tate module
of \(B_8^5\) has no Galois-stable rational line:
\cite[\S1, Corollary 1]{Ruppert1998} would otherwise give a
nonzero CM abelian subvariety. For \(A_2\), the CM type
\(\{1,2\}\) has trivial stabilizer in
\((\mathbf Z/5)^\times\). Hence the four characters obtained
by splitting \(F\otimes\overline{\mathbf Q}_\ell\)
remain pairwise distinct on every open global Galois subgroup.
Here is the needed labelled argument. Equality of two characters
on such a subgroup implies equality of their Hodge--Tate weights
at every embedding of a finite field of definition into
\(\overline{\mathbf Q}_\ell\). The labelled weights are the
translates of the indicator of \(\{1,2\}\); its trivial
stabilizer forces the two embeddings to be identical.
The comparison and invariance under finite extension are
\cite[Theorems 2.2.3 and 2.4.6, pp.13 and 20]{fghBrinonConrad}.

If \(\ell\not\equiv1\pmod5\), every local field factor of
\(F\otimes\mathbf Q_\ell\) has degree greater than one
(including the totally ramified factor at five). A stable
\(\mathbf Q_\ell\)-line in \(V_\ell(A_2^5)\) would, after
projection to a nonzero factor and extension of scalars, have
the same character in every embedding in that local orbit.
The orbit has size greater than one, contradicting the
distinctness just proved. Thus no such line exists.

Now a noncuspidal normalized procyclic subgroup in \(X\)
has, by A.11, a nonzero proper-homology vector spanning a
line under an open arithmetic subgroup. Its projections to
\(A_2^5\) and \(B_8^5\) vanish by the preceding paragraphs.
It therefore projects nontrivially to the original elliptic
factor. The actual map \([5]q:X\to U\) is nonzero on precisely
that factor. Its image procyclic subgroup is nontrivial and
still has open arithmetic normalizer. The original normalizer
theorem forces that image to be cuspidal, with zero proper
homology, a contradiction. This proves the claimed normalizer
classification. Conversely every nontrivial cusp-inertia subgroup
has the required open arithmetic normalizer after making its
cusp rational.

At split primes the remaining possible proper line lies in
\(A_2^5\). The actual genus-two quotient in (A.54.7), with
its three cusps \(z=0\) and infinity, has a degree-five etale
map to the tripod via
\(t=(x+\sqrt\lambda)/(2\sqrt\lambda)\): its three missing
values are \(0,1,\infty\), and after scaling its equation is
\(w^5=t(t-1)\). The deck translates of the actual map
\(C\to D_\lambda/\sigma\) jointly detect the whole
\(A_2^5\) homology: their image is the \(H_5\)-span of the
\(H\)-invariant line in the simple module, tensored with the
one-dimensional \(T=-1\) multiplicity space. Thus a normalizer
theorem for this specific three-punctured genus-two target
would settle the remaining split-prime source question.
For \(\ell\ne5\) its degree-five map to the tripod does not
make its pro-\(\ell\) group an open subgroup of the tripod
group, so no split-prime conclusion follows from that map alone.
The executed exact certificate is retained with this proof.
\end{proof}
